%% kapa: exam=1 round=33 year=2022 period=2 form=A subject=accounting subject_ko=회계학 seq=12 number=52
%% source: Q-Net 감정평가사 기출문제 (cst00309) artlSeq=5209997 file=raw/1008/5209997/2022년도 제33회 감정평가사 1차시험 2교시(감정평가 관계법규, 회계학).pdf
%% conversion: kordoc 4.24.51 -> tools/gen1.py (auto); review=proofread-sonnet; flags=none
\begin{question}{52}
(주)감평은 20×1년 1월 1일 다음과 같은 조건의 전환사채를 액면발행하였다.
\begin{framedtext}
○ 액면금액: ￦1,000,000 \par
○ 표시이자율: 연 6\% \par
○ 일반사채 시장이자율: 연 10\% \par
○ 이자지급일: 매년 말 \par
○ 만기상환일: 20×3년 12월 31일 \par
\end{framedtext}
동 전환사채는 전환권을 행사하지 않을 경우 만기상환일에 액면금액의 106.49\%를 일시 상환하는 조건이다. 전환청구가 없었다고 할 때, (주)감평이 동 전환사채와 관련하여 3년(20×1년 1월 1일\textasciitilde 20×3년 12월 31일)간 인식할 이자비용 총액은? (단, 단수차이로 인한 오차가 있다면 가장 근사치를 선택한다.)
\par{\small\centering\begin{tabular}{|c|c|c|c|c|}\hline
\multirow{2}{*}{기간} & \multicolumn{2}{c|}{단일금액 ￦1의 현재가치} & \multicolumn{2}{c|}{정상연금 ￦1의 현재가치} \\ \cline{2-5}
 & 6\% & 10\% & 6\% & 10\% \\ \hline
3 & 0.83962 & 0.75131 & 2.67301 & 2.48685 \\ \hline
\end{tabular}\par}
\begin{choices}
  \item ￦50,719
  \item ￦115,619
  \item ￦244,900
  \item ￦295,619
  \item ￦344,619
\end{choices}
\end{question}
\answer{4}
